\documentclass[12pt,a4paper]{report}

\usepackage[utf8]{inputenc}
\usepackage[T1]{fontenc}
\usepackage[english]{babel}

\usepackage{geometry}
\geometry{
    a4paper,
    left=30mm, 
    right=20mm,
    top=25mm,
    bottom=25mm
}
\usepackage{setspace}
\onehalfspacing 
\usepackage{lmodern} 

\usepackage{amsmath, amssymb, amsthm}
\usepackage{algorithm}
\usepackage{algpseudocode}

\usepackage{graphicx}
\usepackage{hyperref}
\hypersetup{
    colorlinks=true,
    linkcolor=blue,
    filecolor=magenta,      
    urlcolor=cyan,
    citecolor=blue,
    pdftitle={Longest Increasing Subsequence},
    pdfauthor={Your Name}
}

\usepackage[backend=biber, style=ieee]{biblatex}
\addbibresource{bibliography.bib} 

\newtheorem{theorem}{Theorem}[chapter]
\newtheorem{lemma}[theorem]{Lemma}
\newtheorem{definition}{Definition}[chapter]

\title{\textbf{Longest Increasing Subsequence Problem in Small Space, revisited\\}}
\author{Mikołaj Owczarz \\ \vspace{0.5cm} \\ 
University of Wrocław}
\date{April 2026}

\begin{document}

\maketitle

\begin{abstract}
ABSTRACT ABSTRACT ABSTRACT ABSTRACT ABSTRACT ABSTRACT

\end{abstract}

\tableofcontents
\newpage

\section{Introduction}
Problem of finding a longest increasing subsequence is a classic problem in Computer Science.
Having wide variety of practical applications including (...) it is a heavy researched topic with many 
variations and different algorithms developed such as ... .

In this paper we focus on solving (both obtaining length and computing the sequence as a whole) LIS problem using small space.
Based on Kiyomi ... we show proof that presented there Patience Sorting computes length of LIS in $O(max(nlogn, n^2/s))$, and 
we provide alternative approach to using their algorithm in D&C to obtain said sequence in $O(max(n*logn*loglogs, n^2/s))$ time.



\subsection{Preliminaries}
Basic definicje, jaki model uzywamy
\subsubsection{LIS problem}
\subsubsection{Classic algorithm} to nie wiem czy trzeba 
\subsubsection{Patience sorting} to wsumie chyba nie tu
\subsection{Previous results}
Co oni zrobili, jak uzyli patience sortingu
\subsubsection{Time complexity for length computation}
\subsubsection{Time complexity for LIS computation}
\subsection{New results}
\subsubsection{Length computation}
In this paper we provide a tighter analysis of Patience Sorting algorithm, the same as one presented by Kiyomi's and ... in [].
We show that this algorithm actually runs in $O(max(n*logn, n^2/s))$ which is tighter than previous $O(n^2logn/s)$ presented in [].
\subsubsection{LIS computation}
For retrieving LIS sequence we provide an alternative Divide & Conquer algorithm, which combined with previous observation runs in 
$O(max(n*logn*loglogs, n^2/s))$ time, improvement on earlier result of $O(n^2/s*log^2(s))$.

\section{Length Computation Complexity Analysis}
\subsection{Patience sorting}
\subsection{Time complexity analysis of patience sorting with parametrised size}

\section{LIS Computation Algorithm and Analysis}

\section{Summary}
\subsection{Open questions}
jakies summary
\printbibliography

\end{document}